\documentclass[10pt,a4paper]{amsart}
\usepackage[margin=19mm]{geometry}
\usepackage{amsmath,amssymb,mathtools}
\usepackage[scr=rsfs,cal=euler]{mathalfa}
\usepackage{xfrac}
\usepackage[unicode,hidelinks,bookmarks=false]{hyperref}
\newcommand{\ie}{i.e.\ }
\newcommand{\cf}{cf.\ }
\newcommand{\resp}{respectively\ }
\newcommand{\Subsection}{\subsection}
\providecommand{\nobreakdash}{\nobreak\hbox{-}}
\setlength{\emergencystretch}{2em}
\multlinegap0pt
\usepackage{amssymb}
\usepackage{mathtools}
\newtheorem{lemma}{Lemma}
\title{The exact generalized Tur\'an number for \(C_6\) in \(C_8\)-free graphs}
\author{Zian Chen, Jinghua Deng}
\date{Source: arXiv:2607.03856v1 (CC BY 4.0)}
\begin{document}
\maketitle

\noindent\emph{Source and licence.} The exact generalized Tur\'an number for \(C_6\) in \(C_8\)-free graphs. Version arXiv:2607.03856v1 (CC BY 4.0), distributed by arXiv under the Creative Commons Attribution 4.0 International licence, http://creativecommons.org/licenses/by/4.0/, as recorded in the arXiv licence metadata for this exact version and shown on the abstract page https://arxiv.org/abs/2607.03856v1. Full licence notice: \texttt{licenses/arxiv-2607.03856-CC-BY-4.0.txt}.\par\smallskip

\noindent\textit{Editorial scope.} Counted unit: the article's lemma lem:one-two-fat (source0.tex lines 342--351). The article's complete original proof of it is delivered with it (source0.tex lines 353--429, one \texttt{proof} environment of 77 lines). No mathematics is written, repaired or re-derived: the statement and the proof are the article's own lines, in the article's own notation, and no retained line is substituted or re-flowed.  Retained uncounted as the source-local context the proof needs: lines 243--245; lines 314--319.  Nothing was omitted from any retained span, so the package carries no omission record.  No retained line contains a citation, so the wrapper carries no bibliography and no bibliography entry.  No retained line contains a figure environment, an image-inclusion command or drawing code, so no image is delivered and the package declares no asset dependency.  Editorial formatting only, in addition: an AMS \texttt{amsart} preamble that restates the article's own declarations and macros used by the retained lines, namely the environments \texttt{lemma} on separate counters, together with the standard packages those lines need.  The retained environments print in this document as Lemma 1 in order of appearance, whereas the article prints the same items with the numbering of its own full document.  The complete native source of the article as distributed in the arXiv e-print for this exact version is bundled unchanged as \texttt{sources/204428/source0.tex}.
\begin{equation}\label{eq:charge}
  F^{\#}(S)\le d(u,v)d(v,w)d(w,u)\le\mu(S)^3.
\end{equation}

\begin{lemma}\label{lem:linear-budget}
For every fixed \(\rho>0\),
\[
  \sum_{e\in E(\tilde{G})}d_G(e)\le3n+O_\rho(1).
\]
\end{lemma}

\begin{lemma}\label{lem:one-two-fat}
For fixed \(K\ge10\), the total contribution of 2-fat charged triples is
\[
  O(\rho n^3)+O_\rho(n^2),
\]
and the total contribution of 1-fat charged triples is
\[
  O_K(\rho n^3)+O_{K,\rho}(n^2).
\]
\end{lemma}

\begin{proof}
We first treat 2-fat triples. Let \(T=\{v_1,v_2,v_3\}\) be 2-fat, and suppose, after relabeling, that
  \(v_1v_2,v_1v_3\in E(\tilde{G})\), and \(v_2v_3\notin E(\tilde{G})\).

By the definition of \(\tilde{G}\), the last condition gives \(d(v_2,v_3)<\rho n\). Therefore, by
\eqref{eq:charge},
\[
  F^{\#}(T)
  \le d(v_1,v_2)d(v_1,v_3)d(v_2,v_3)
  \le \rho n\,d(v_1,v_2)d(v_1,v_3).
\]
For each 2-fat triple \(T\), let \(e_1(T),e_2(T)\) be its two fat pairs. Summing the previous bound
and then enlarging the sum to all ordered pairs of edges of \(\tilde{G}\), we obtain
\[
  \sum_{T\text{ 2-fat}}F^{\#}(T)
  \le \rho n\sum_{T\text{ 2-fat}}d(e_1(T))d(e_2(T))
  \le \rho n\sum_{e_1,e_2\in E(\tilde{G})}d(e_1)d(e_2)
  =\rho n\left(\sum_{e\in E(\tilde{G})}d(e)\right)^2.
\]
Lemma~\ref{lem:linear-budget} gives
\[
  \sum_{T\text{ 2-fat}}F^{\#}(T)
  \le \rho n(3n+O_\rho(1))^2
  =O(\rho n^3)+O_\rho(n^2).
\]


We next treat 1-fat triples. Fix an edge \(v_1v_2\in E(\tilde{G})\), and for
\(\omega_0\in V(G)\setminus\{v_1,v_2\}\), write
\[
  p_{\omega_0}=d(v_1,\omega_0),\qquad q_{\omega_0}=d(v_2,\omega_0).
\]
We claim that
\begin{equation}\label{eq:pq}
  \sum_{\substack{\omega_0\in V(G)\setminus\{v_1,v_2\}\\p_{\omega_0},q_{\omega_0}<\rho n}}
  p_{\omega_0}q_{\omega_0}
  =O_K(\rho n^2).
\end{equation}
Indeed, vertices \(\omega_0\) with \(\min\{p_{\omega_0},q_{\omega_0}\}<K\) contribute at most \(K\rho n\)
each, since \(p_{\omega_0},q_{\omega_0}<\rho n\). Thus their total contribution is at most
\(K\rho n^2\).

There is at most one vertex \(\omega_0\) with \(p_{\omega_0},q_{\omega_0}\ge K\). Indeed, if two
distinct vertices \(\omega_1,\omega_2\) satisfied
$p_{\omega_1},q_{\omega_1},p_{\omega_2},q_{\omega_2}\ge K$,
then, since \(K\ge10\), we could choose vertices
\[
  \omega_3\in N(v_1,\omega_1),\quad
  \omega_4\in N(\omega_1,v_2),\quad
  \omega_5\in N(v_2,\omega_2),\quad
  \omega_6\in N(\omega_2,v_1)
\]
successively so that they are distinct and avoid the already named vertices
\(\{v_1,v_2,\omega_1,\omega_2\}\). Indeed, at each step at most seven vertices are forbidden,
whereas each relevant common neighborhood has size at least \(K\).
Then
  $v_1$-$\omega_3$-$\omega_1$-$\omega_4$-$v_2$-$\omega_5$-$\omega_2$-$\omega_6$-$v_1$
is a \(C_8\), a contradiction. Hence at most one remaining vertex exists, and its contribution is at
most \((\rho n)^2\le\rho n^2\). This proves \eqref{eq:pq}.

For a 1-fat triple whose unique fat pair is \(v_1v_2\), the third vertex \(\omega_0\) satisfies
\(p_{\omega_0},q_{\omega_0}<\rho n\), and
\[
  F^{\#}(\{v_1,v_2,\omega_0\})\le d(v_1,v_2)p_{\omega_0}q_{\omega_0}.
\]
Therefore, using \eqref{eq:pq} for each fixed fat edge and then Lemma~\ref{lem:linear-budget},
\[
  \sum_{T\text{ 1-fat}}F^{\#}(T)
  \le
  \sum_{v_1v_2\in E(\tilde{G})}d(v_1,v_2)
  \sum_{\substack{\omega_0\in V(G)\setminus\{v_1,v_2\}\\p_{\omega_0},q_{\omega_0}<\rho n}}
  p_{\omega_0}q_{\omega_0}
  \le O_K(\rho n^2)\sum_{e\in E(\tilde{G})}d(e)
  =O_K(\rho n^3)+O_{K,\rho}(n^2).
\]
This proves the 1-fat estimate.
\end{proof}

\end{document}
